\section{AMC 12 12B 2011}
        \begin{enumerate}
        	\item (\textit{AMC 12 12B 2011}) 
What is 

 $\frac{2+4+6}{1+3+5}-\frac{1+3+5}{2+4+6}?$
$\textbf{(A)}\ -1 \qquad \textbf{(B)}\ \frac{5}{36} \qquad \textbf{(C)}\ \frac{7}{12} \qquad \textbf{(D)}\ \frac{147}{60} \qquad \textbf{(E)}\ \frac{43}{3}$

 \textbf{Answer: }C



\textbf{Solution 1}

Add up the numbers in each fraction to get $\frac{12}{9} - \frac{9}{12}$, which equals $\frac{4}{3} - \frac{3}{4}$. Doing the subtraction yields $\boxed{\frac{7}{12}\   \textbf{(C)}}$




\textbf{Solution 2}

Notice that the numerators and denominators of each expression are 3-term arithmetic series. The sum of an arithmetic series is the middle term multiplied with the number of terms. Since each of the arithmetic series have the same number of terms, we can replace the fractions with the middle terms, which gives us $\frac{4}{3} - \frac{3}{4} = \frac{7}{12}$



	\item (\textit{AMC 12 12B 2011}) Josanna's test scores to date are $90$, $80$, $70$, $60$, and $85$.  Her goal is to raise her test average at least $3$ points with her next test.  What is the minimum test score she would need to accomplish this goal?


 $\textbf{(A)}\ 80 \qquad \textbf{(B)}\ 82 \qquad \textbf{(C)}\ 85 \qquad \textbf{(D)}\ 90 \qquad \textbf{(E)}\ 95$

 \textbf{Answer: }E



\textbf{Solution}

Take the average of her current test scores, which is
 \[\frac{90+80+70+60+85}{5} = \frac{385}{5} = 77\]

 This means that she wants her test average after the sixth test to be $80.$ Let $x$ be the score that Josanna receives on her sixth test. Thus, our equation is

 
 \[\frac{90+80+70+60+85+x}{6} = 80\]

 \[385+x = 480\]

 \[x = \boxed{95 \textbf{(E)}}\]



	\item (\textit{AMC 12 12B 2011}) LeRoy and Bernardo went on a week-long trip together and agreed to share the costs equally.  Over the week, each of them paid for various joint expenses such as gasoline and car rental.  At the end of the trip it turned out that LeRoy had paid $A$ dollars and Bernardo had paid $B$ dollars, where $A<B$.  How many dollars must LeRoy give to Bernardo so that they share the costs equally?


 $\textbf{(A)}\ \frac{A+B}{2} \qquad \textbf{(B)}\ \frac{A-B}{2} \qquad \textbf{(C)}\ \frac{B-A}{2} \qquad \textbf{(D)}\ B-A \qquad \textbf{(E)}\ A+B$

 \textbf{Answer: }C



\textbf{Solution}

The total amount of money that was spent during the trip was
 \[A + B\]So each person should pay
 \[\frac{A+B}{2}\]if they were to share the costs equally. Because LeRoy has already paid $A$ dollars of his part, he still has to pay
 \[\frac{A+B}{2} - A =\]
 \[= \boxed{\frac{B-A}{2} \textbf{(C)}}\]



	\item (\textit{AMC 12 12B 2011}) In multiplying two positive integers $a$ and $b$, Ron reversed the digits of the two-digit number $a$.  His erroneous product was 161.  What is the correct value of the product of $a$ and $b$?


 $\textbf{(A)}\ 116 \qquad \textbf{(B)}\ 161 \qquad \textbf{(C)}\ 204 \qquad \textbf{(D)}\ 214 \qquad \textbf{(E)}\ 224$

 \textbf{Answer: }E



\textbf{Solution}

Taking the prime factorization of $161$ reveals that it is equal to $23*7.$ Therefore, the only ways to represent $161$ as a product of two positive integers is $161*1$ and $23*7.$ Because neither $161$ nor $1$ is a two-digit number, we know that $a$ and $b$ are $23$ and $7.$ Because $23$ is a two-digit number, we know that a, with its two digits reversed, gives $23.$ Therefore, $a = 32$ and $b = 7.$ Multiplying our two correct values of $a$ and $b$ yields

 
 \[a*b = 32*7 =\]

 \[= \boxed{224 \textbf{(E)}}\]



	\item (\textit{AMC 12 12B 2011}) Let $N$ be the second smallest positive integer that is divisible by every positive integer less than $7$.  What is the sum of the digits of $N$?


 $\textbf{(A)}\ 3 \qquad \textbf{(B)}\ 4 \qquad \textbf{(C)}\ 5 \qquad \textbf{(D)}\ 6 \qquad \textbf{(E)}\ 9$

 \textbf{Answer: }A



\textbf{Solution}

$N$ must be divisible by every positive integer less than $7$, or $1, 2, 3, 4, 5,$ and $6$. Each number that is divisible by each of these is a multiple of their least common multiple. $LCM(1,2,3,4,5,6)=60$, so each number divisible by these is a multiple of $60$. The smallest multiple of $60$ is $60$, so the second smallest multiple of $60$ is $2\times60=120$. Therefore, the sum of the digits of $N$ is $1+2+0=\boxed{3\ \textbf{(A)}}$



	\item (\textit{AMC 12 12B 2011}) Two tangents to a circle are drawn from a point $A$.  The points of contact $B$ and $C$ divide the circle into arcs with lengths in the ratio $2 : 3$.  What is the degree measure of $\angle{BAC}$?


 $\textbf{(A)}\ 24 \qquad \textbf{(B)}\ 30 \qquad \textbf{(C)}\ 36 \qquad \textbf{(D)}\ 48 \qquad \textbf{(E)}\ 60$

 \textbf{Answer: }C



\textbf{Solution}

In order to solve this problem, use of the tangent-tangent intersection theorem (Angle of intersection between two tangents dividing a circle into arc length A and arc length B = 1/2 (Arc A° - Arc B°).

 In order to utilize this theorem, the degree measures of the arcs must be found.  First, set A (Arc length A) equal to 3d, and B (Arc length B) equal to 2d.  

 Setting 3d+2d = 360° will find d = 72°, and so therefore Arc length A in degrees will equal 216° and arc length B will equal 144°.

 Finally, simply plug the two arc lengths into the tangent-tangent intersection theorem, and the answer:

 1/2 (216°-144°) = 1/2 (72°) $=\boxed{36 \textbf{(C)}}.$




\textbf{Solution 2}

Let arc A be 3x and arc B be 2x. Then the angle formed by the tangents is $\frac{3x-2x}2 = \frac {x} 2$ by the arc length formula. Also note that $3x + 2x = 360$, which simplifies to $x= 72.$ Hence the angle formed by the tangents is equal to $\boxed{36 \textbf{(C)}}$.

 



\textbf{Solution 3}

Let the center of the circle be $O$. The radii extending from $O$ to the points of tangency $B$ and $C$ are both perpendicular to lines $AB$ and $AC$. Thus $\angle ABO = \angle ACO = 90^{\circ}$, and quadrilateral $ABOC$ is cyclic, since the opposite angles add to $180^{\circ}$. 

 Since the ratio of the arc lengths is $2 : 3$, the shorter arc (the arc cut off by $\angle BOC$) is $144^{\circ}$. 

 $\angle BOC$ and $\angle BAC$ must add up to $180^{\circ}$ since quadrilateral $ABOC$ is cyclic, so $\angle BAC = 180 - 144 = \boxed{\textbf{(C) } 36}$.

 -Solution by Joeya

 


	\item (\textit{AMC 12 12B 2011}) Let $x$ and $y$ be two-digit positive integers with mean $60$.  What is the maximum value of the ratio $\frac{x}{y}$?


 $\textbf{(A)}\ 3 \qquad \textbf{(B)}\ \frac{33}{7} \qquad \textbf{(C)}\ \frac{39}{7} \qquad \textbf{(D)}\ 9 \qquad \textbf{(E)}\ \frac{99}{10}$

 \textbf{Answer: }B



\textbf{Solution}

If $x$ and $y$ have a mean of $60$, then $\frac{x+y}{2}=60$ and $x+y=120$. To maximize $\frac{x}{y}$, we need to maximize $x$ and minimize $y$. Since they are both two-digit positive integers, the maximum of $x$ is $99$ which gives $y=21$. $y$ cannot be decreased because doing so would increase $x$, so this gives the maximum value of $\frac{x}{y}$, which is $\frac{99}{21}=\boxed{\frac{33}{7}\ \textbf{(B)}}$



	\item (\textit{AMC 12 12B 2011}) Keiko walks once around a track at exactly the same constant speed every day. The sides of the track are straight, and the ends are semicircles. The track has width $6$ meters, and it takes her $36$ seconds longer to walk around the outside edge of the track than around the inside edge. What is Keiko's speed in meters per second?


 $\textbf{(A)}\ \frac{\pi}{3} \qquad \textbf{(B)}\ \frac{2\pi}{3} \qquad \textbf{(C)}\ \pi \qquad \textbf{(D)}\ \frac{4\pi}{3} \qquad \textbf{(E)}\ \frac{5\pi}{3}$

 \textbf{Answer: }A



\textbf{Solution}

To find Keiko's speed, all we need to find is the difference between the distance around the inside edge of the track and the distance around the outside edge of the track, and divide it by the difference in the time it takes her for each distance.  We are given the difference in time, so all we need to find is the difference between the distances.

 The track is divided into lengths and curves.  The lengths of the track will exhibit no difference in distance between the inside and outside edges, so we only need to concern ourselves with the curves.

 The curves of the track are semicircles, but since there are two of them, we can consider both of the at the same time by treating them as a single circle.  We need to find the difference in the circumferences of the inside and outside edges of the circle.

 The formula for the circumference of a circle is $C = 2 * \pi * r$ where $r$ is the radius of the circle.

 Let's define the circumference of the inside circle as $C_1$ and the circumference of the outside circle as $C_2$.

 If the radius of the inside circle ($r_1$) is $n$, then given the thickness of the track is 6 meters, the radius of the outside circle ($r_2$) is $n + 6$.

 Using this, the difference in the circumferences is:

 $C_2 - C_1 = 2 * \pi * (r_2 - r_1) = 2 * \pi * (n + 6 - n) = 2 * \pi * 6 = 12\pi$
$12\pi$ is the difference between the inside and outside lengths of the track.  Divided by the time differential, we get:

 $12\pi \div 36 = \boxed {\textbf{(A)}\ \frac{\pi}{3}}$



	\item (\textit{AMC 12 12B 2011}) Two real numbers are selected independently and at random from the interval $[-20,10]$.  What is the probability that the product of those numbers is greater than zero?


 $\textbf{(A)}\ \frac{1}{9} \qquad \textbf{(B)}\ \frac{1}{3} \qquad \textbf{(C)}\ \frac{4}{9} \qquad \textbf{(D)}\ \frac{5}{9} \qquad \textbf{(E)}\ \frac{2}{3}$

 \textbf{Answer: }D



\textbf{Solution}

For the product to be greater than zero, we must have either both numbers negative or both positive. 

 Both numbers are negative with a $\frac{2}{3}*\frac{2}{3}=\frac{4}{9}$ chance.

 Both numbers are positive with a $\frac{1}{3}*\frac{1}{3}=\frac{1}{9}$ chance.

 Therefore, the total probability is $\frac{4}{9}+\frac{1}{9}=\frac{5}{9}$ and we are done. $\boxed{D}$



	\item (\textit{AMC 12 12B 2011}) Rectangle $ABCD$ has $AB=6$ and $BC=3$. Point $M$ is chosen on side $AB$ so that $\angle AMD=\angle CMD$. What is the degree measure of $\angle AMD$?


 $\textbf{(A)}\ 15 \qquad \textbf{(B)}\ 30 \qquad \textbf{(C)}\ 45 \qquad \textbf{(D)}\ 60 \qquad \textbf{(E)}\ 75$

 \textbf{Answer: }E



\textbf{Solution}

Since $AB \parallel CD$, $\angle AMD = \angle CDM$, so $\angle AMD = \angle CMD = \angle CDM$, so $\bigtriangleup CMD$ is isosceles, and hence $CM=CD=6$.  Therefore, $\triangle BMC$ is a 30-60-90 triangle with $\angle BMC = 30^\circ$.  Therefore $\angle AMD=\boxed{\mathrm{(E)}\ 75^\circ}$



	\item (\textit{AMC 12 12B 2011}) A frog located at $(x,y)$, with both $x$ and $y$ integers, makes successive jumps of length $5$ and always lands on points with integer coordinates. Suppose that the frog starts at $(0,0)$ and ends at $(1,0)$. What is the smallest possible number of jumps the frog makes?


 $\textbf{(A)}\ 2 \qquad \textbf{(B)}\ 3 \qquad \textbf{(C)}\ 4 \qquad \textbf{(D)}\ 5 \qquad \textbf{(E)}\ 6$

 \textbf{Answer: }B



\textbf{Solution}

Since the frog always jumps in length $5$ and lands on a lattice point, the sum of its coordinates must change either by $5$ (by jumping parallel to the x- or y-axis), or by $3$ or $4$ (3-4-5 right triangle).

 Because either $1$, $5$, or $7$ is always the change of the sum of the coordinates, the sum of the coordinates will always change from odd to even or vice versa. Thus, it can't go from $(0,0)$ to $(1,0)$ in an even number of moves. Therefore, the frog cannot reach $(1,0)$ in two moves.  

 However, a path is possible in 3 moves:  from $(0,0)$ to $(3,4)$ to $(6,0)$ to $(1,0)$.

 Thus, the answer is  $= \boxed{3 \textbf{}}$.

 


	\item (\textit{AMC 12 12B 2011}) A dart board is a regular octagon divided into regions as shown below. Suppose that a dart thrown at the board is equally likely to land anywhere on the board. What is the probability that the dart lands within the center square?


 
                    \begin{figure}[H]
                        \centering
                        \includegraphics[width=0.5\textwidth]{Problem/AMC_12_12B_2011_12_Problem_Figure_1.png}
                    \end{figure}
                    
$\textbf{(A)}\ \frac{\sqrt{2} - 1}{2} \qquad \textbf{(B)}\ \frac{1}{4} \qquad \textbf{(C)}\ \frac{2 - \sqrt{2}}{2} \qquad \textbf{(D)}\ \frac{\sqrt{2}}{4} \qquad \textbf{(E)}\ 2 - \sqrt{2}$

 \textbf{Answer: }A



\textbf{Solution}

Let's assume that the side length of the octagon is $x$. The area of the center square is just $x^2$. The triangles are all $45-45-90$ triangles, with a side length ratio of $1:1:\sqrt{2}$. The area of each of the $4$ identical triangles is $\left(\dfrac{x}{\sqrt{2}}\right)^2\times\dfrac{1}{2}=\dfrac{x^2}{4}$, so the total area of all of the triangles is also $x^2$. Now, we must find the area of all of the 4 identical rectangles. One of the side lengths is $x$ and the other side length is $\dfrac{x}{\sqrt{2}}=\dfrac{x\sqrt{2}}{2}$, so the area of all of the rectangles is $2x^2\sqrt{2}$. The ratio of the area of the square to the area of the octagon is $\dfrac{x^2}{2x^2+2x^2\sqrt{2}}$. Cancelling $x^2$ from the fraction, the ratio becomes $\dfrac{1}{2\sqrt2+2}$. Multiplying the numerator and the denominator each by $2\sqrt{2}-2$ will cancel out the radical, so the fraction is now $\dfrac{1}{2\sqrt2+2}\times\dfrac{2\sqrt{2}-2}{2\sqrt{2}-2}=\dfrac{2\sqrt{2}-2}{4}=\boxed{\mathrm{(A)}\ \dfrac{\sqrt{2}-1}{2}}$




\textbf{Solution}

Instead of doing mindless algebra, let's set the side lengths of the square to $1$.

 After careful inspection of the regular octagon, we find that all of the sides of the regular octagon must be $1$.

 This gives the total area to be $4 \cdot \text{triangle area}$ + $4 \cdot \text{rectangle area}$ + $\text{square area}$ = $4(\dfrac{1}{4}) + 4(\sqrt{\dfrac{1}{2}}) + 1$, which simplifies to $2(1+\sqrt{2})$.

 The percent area of the center square is the area of the square divided by the total area of the regular octagon. This gives $\dfrac{1}{2(1+\sqrt{2})}.$ Multiplying by the conjugate gives answer choice $\boxed{(A).}$



	\item (\textit{AMC 12 12B 2011}) Brian writes down four integers $w > x > y > z$ whose sum is $44$. The pairwise positive differences of these numbers are $1, 3, 4, 5, 6$ and $9$. What is the sum of the possible values of $w$?


 $\textbf{(A)}\ 16 \qquad \textbf{(B)}\ 31 \qquad \textbf{(C)}\ 48 \qquad \textbf{(D)}\ 62 \qquad \textbf{(E)}\ 93$

 \textbf{Answer: }B



\textbf{Solution}

Assume that $y-z=a, x-y=b, w-x=c.$ $w-z$ results in the greatest pairwise difference, and thus it is $9$. This means $a+b+c=9$. $a,b,c$ must be in the set ${1,3,4,5,6}$.The only way for 3 numbers in the set to add up to 9 is if they are $1,3,5$. $a+b$, and $b+c$ then must be the remaining two numbers which are $4$ and $6$.The ordering of $(a,b,c)$ must be either $(3,1,5)$ or $(5,1,3)$.

 

 \begin{align*} z + (z + a) + (z + (a + b)) + (z + (a + b + c)) &= 4z + a + (a + b) + 9\\ 4z + a + (a + b) + 9 &= 44\\ \text{if} \hspace{1cm} a &= 3 \\ a + b &= 4\\ 4z &= 44 - 9 - 3 - 4\\ z &= 7\\ w &= 16\\ \end{align*}


 \begin{align*} \text{if} \hspace{1cm} a &= 5\\ a + b &= 6\\ 4z &= 44 - 9 - 5 - 6\\ z &= 6\\ w &= 15\\ \end{align*}


 
The sum of the two w's is $15+16=31$ $\boxed{B}$




\textbf{Solution 2}

Let the four numbers be $z$, $z+a$, $z+b$, and $z+c$. We know that $c$ must be $9$ because that's the greatest difference. So we have $z$, $z+a$, $z+b$, and $z+9$. The 6 possible differences are $a$, $b$, $9$, $b-a$, $9-a$, and $9-b$. We are given that the differences are 1, 3, 4, 5, 6, 9. $a$ and $9-a$ and $b$ and $9-b$ add to 9 which means they have to be 4, 5 and 3,6 or vice versa. Which leaves $1$. That means $b-a$ has to equal $1$. So a and b have to be 3,4, 4,5, or 5,6. For 3,4, we have $z$, $z+3$, $z+4$, and $z+9$. $4z+16=44$. $4z=28$. $z=7$. $w=7+9=16$. Now for 4,5, notice that it doesn't work. The differences are 4, 5, 9, 1, 4, 5. We are missing 6 and 3. For 5,6, it's $z$, $z+5$, $z+6$, and $z+9$. Check that the differences work; they do. We have $4z+20=44$. $4z=24$. $z=6$. $w=6+9=15$. Therefore our answer is $15+16=\boxed{31}$.~MC413551

 


	\item (\textit{AMC 12 12B 2011}) A segment through the focus $F$ of a parabola with vertex $V$ is perpendicular to $\overline{FV}$ and intersects the parabola in points $A$ and $B$. What is $\cos\left(\angle AVB\right)$?


 $\textbf{(A)}\ -\frac{3\sqrt{5}}{7} \qquad \textbf{(B)}\ -\frac{2\sqrt{5}}{5} \qquad \textbf{(C)}\ -\frac{4}{5} \qquad \textbf{(D)}\ -\frac{3}{5} \qquad \textbf{(E)}\ -\frac{1}{2}$

 \textbf{Answer: }D



\textbf{Solution  1}

Name the directrix of the parabola $l$. Define $d(X,k)$ to be the distance between a point $X$ and a line $k$.

 Now we remember the geometric definition of a parabola: given any line $l$ (called the directrix) and any point $F$ (called the focus), the \textit{parabola} corresponding to the given directrix and focus is the locus of the points that are equidistant from $F$ and $l$. Therefore $FV=d(V,l)$. Let this distance be $d$. Now note that $d(F,l)=2d$, so $d(A,l)=d(B,l)=2d$. Therefore $AF=BF=2d$. We now use the \href{https://artofproblemsolving.com/wiki/index.php?title=Pythagorean\_Theorem}{Pythagorean Theorem} on triangle $AFV$; $AV=\sqrt{AF^2+FV^2}=d\sqrt{5}$. Similarly, $BV=d\sqrt{5}$. We now use the \href{https://artofproblemsolving.com/wiki/index.php?title=Law\_of\_Cosines}{Law of Cosines}:

 
 \[AB^2=AV^2+VB^2-2AV\cdot VB\cos{\angle AVB}\Rightarrow 16d^2=10d^2-10d^2\cos{\angle AVB}\]

 \[\Rightarrow \cos{\angle AVB}=-\frac{3}{5}\]

 This shows that the answer is $\boxed{\textbf{(D)}}$.

 


 



\textbf{Solution 2}

WLOG we can assume that the parabola is $y=x^2$. Therefore $V = (0,0)$ and $F = (0,\frac{1}{4})$. Also $A = (-\frac{1}{2},\frac{1}{4})$ and $B = (\frac{1}{2},\frac{1}{4})$.

 $AB = 1$ and $AV = VB = \sqrt{(\frac{1}{2})^2+(\frac{1}{4})^2} = \frac{\sqrt{5}}{4}$ by the pythagorean theorem.

 Now using the law of cosines on $\triangle AVB$ we have: 

 $AB^2 = 2AV^2-2AV^2\cos(\angle AVB) = 2AV^2(1-\cos(\angle AVB))$
$1 = \frac{5}{8} (1-\cos(\angle AVB))$

 Thus, 
 \[\cos(\angle AVB) = \boxed{\textbf{(D)} -\frac{3}{5}}.\]

 (solution by mihirb)

 



\textbf{Solution 3}

After assuming that the parabola is $x^2$,find the points A and B, which are +/- 1,2,1/4.Now treat them as vectors,take the dot product,then find the magnitudes and multiply them.A well known definition of the dot product says that the quotient of the two is the cosine of the angle between them.This will give you D.

 



\textbf{Solution 4}

As we know, an expression for the equation of a parabola is $y-k=\frac{1}{4a}(x-h)^2$, where $(h,k)$ is the vertex and $a$ is the distance from the focus to the vertex, here $F$ to $V$. The length of the latus rectum, or $\overline{AB}$ here, is equal to $|4a|$. This means that $AF=BF=2a$ since $F$ is the midpoint of $AB$. Then we can use right Triangle $VAF$ to figure out that $VA=\frac{a\sqrt{5}}{4}$. Now we can use the fact that $\cos{\angle AVB}=\cos{2 \angle AVF}$ and use the double angle formula. This results in $\cos{2 \angle AVF}=2{\cos}^2{\angle AVF}-1$. We can find $\cos{\angle AVF}$ from right triangle $AVF$ using the Pythagorean Theorem, which is $\frac{a}{\sqrt{5}}$. Evaulating the expression, we find that $\cos{2 \angle AVF}=\frac{-3}{5}$.

 -Indefintense

 



\textbf{Solution 5}

Start by drawing a nice diagram. Next, let the directrix of the parabola be $\ell$. We then also let the intersection of line $FV$ and directrix $\ell$ be $X$. We can now let $FV=VX=x$. Now, we draw a segment from $A$ perpendicular to $\ell$ and call the foot of this altitude $Y$. Since $AY=2x$, $AF=2x$. Similarly, $BF=2x$. So, $AV=BV=x\sqrt5$. So, LoC states that 
 \[16x^2=10x^2-10x^2\cos{\angle{AVB}}\] 
 \[\Longrightarrow \cos{\angle{AVB}}=\boxed{-\dfrac{3}{5}}.\]

 -jb2015007

 


	\item (\textit{AMC 12 12B 2011}) How many positive two-digit integers are factors of $2^{24}-1$?


 $\textbf{(A)}\ 4 \qquad \textbf{(B)}\ 8 \qquad \textbf{(C)}\ 10 \qquad \textbf{(D)}\ 12 \qquad \textbf{(E)}\ 14$

 \textbf{Answer: }D



\textbf{Solution}

Repeating \href{https://artofproblemsolving.com/wiki/index.php?title=Difference\_of\_squares}{difference of squares}:

 $2^{24}-1=(2^{12}+1)(2^{6}+1)(2^{3}+1)(2^{3}-1)$
$2^{24}-1=(2^{12}+1)\cdot65\cdot9\cdot7$
$2^{24}-1 = (2^{12} +1) * 5 * 13 * 3^2 * 7$

 The sum of cubes formula gives us: 

 $2^{12}+1=(2^4+1)(2^8-2^4+1)$
$2^{12}+1 = 17\cdot241$

 
A quick check shows $241$ is prime. Thus, the only factors to be concerned about are $3^2\cdot5\cdot7\cdot13\cdot17$, since multiplying by $241$ will make any factor too large.

 Multiplying $17$ by $3$ or $5$ will give a two-digit factor; $17$ itself will also work. The next smallest factor, $7$, gives a three-digit number.  Thus, there are $3$ factors that are multiples of $17$.

 Multiplying $13$ by $3$, $5$, or $7$ will also give a two-digit factor, as well as $13$ itself. Higher numbers will not work, giving $4$ additional factors.

 Multiply $7$ by $3$, $5$, or  $3^2$ for a two-digit factor. There are no more factors to check, as all factors which include $13$ are already counted. Thus, there are an additional $3$ factors.

 Multiply $5$ by $3$ or $3^2$ for a two-digit factor. All higher factors have been counted already, so there are $2$ more factors.

 Thus, the total number of factors is $3+4+3+2=\boxed{\textbf{(D) }12}$

 The 12 two-digit factors are 13, 15, 17, 21, 35, 39, 45, 51, 63, 65, 85, and 91.

 


	\item (\textit{AMC 12 12B 2011}) Rhombus $ABCD$ has side length $2$ and $\angle B = 120^{\circ}$. Region $R$ consists of all points inside of the rhombus that are closer to vertex $B$ than any of the other three vertices. What is the area of $R$?


 $\textbf{(A)}\ \frac{\sqrt{3}}{3} \qquad \textbf{(B)}\ \frac{\sqrt{3}}{2} \qquad \textbf{(C)}\ \frac{2\sqrt{3}}{3} \qquad \textbf{(D)}\ 1 + \frac{\sqrt{3}}{3} \qquad \textbf{(E)}\ 2$

 \textbf{Answer: }C



\textbf{Solution 1}

Suppose that $P$ is a point in the rhombus $ABCD$ and let $\ell_{BC}$ be the \href{https://artofproblemsolving.com/wiki/index.php?title=Perpendicular\_bisector}{perpendicular bisector} of $\overline{BC}$. Then $PB < PC$ if and only if $P$ is on the same side of $\ell_{BC}$ as $B$. The line $\ell_{BC}$ divides the plane into two half-planes; let $S_{BC}$ be the half-plane containing $B$. Let us define similarly $\ell_{BD},S_{BD}$ and $\ell_{BA},S_{BA}$. Then $R$ is equal to $ABCD \cap S_{BC} \cap S_{BD} \cap S_{BA}$. The region turns out to be an irregular pentagon. We can make it easier to find the area of this region by dividing it into four triangles:

 
                    \begin{figure}[H]
                        \centering
                        \includegraphics[width=0.5\textwidth]{Solution/AMC_12_12B_2011_16_Solution_Figure_1.png}
                    \end{figure}
                    Since $\triangle BCD$ and $\triangle BAD$ are equilateral, $\ell_{BC}$ contains $D$, $\ell_{BD}$ contains $A$ and $C$, and $\ell_{BA}$ contains $D$. Then $\triangle BEF \cong \triangle BGF \cong \triangle BGH \cong \triangle BIH$ with $BE = 1$ and $EF = \frac{1}{\sqrt{3}}$ so $[BEF] = \frac{1}{2}\cdot 1 \cdot \frac{\sqrt{3}}{3}$. Multiply this by 4 and it turns out that the pentagon has area $\boxed{(C)\frac{2\sqrt{3}}{3}}$.

 



\textbf{Solution 2}

We follow the steps shown above until we draw pentagon $BIHFE$. We know that rhombus $ABCD$ can be divided into equilateral triangles $\triangle ABD$ and $\triangle CBD$. Using the $30-60-90$ special right triangle rules, we find the height of the equilateral triangles (and the height of the rhombus) to be $\sqrt{3}$. Therefore, the area of $ABCD$ is $2\sqrt{3}$. We now have to take off the areas $\triangle CDA$, $\triangle CEF$, and $\triangle AIH$ to get the desired shape. $\triangle CDA$ is just half of $ABCD$ $(\sqrt {3})$ and $\triangle AIH$ and $\triangle CEF$ are each $\frac{\sqrt {3}}{6}$, for a total area of $2\sqrt {3}-\sqrt {3}-\frac{\sqrt{3}}{6}-\frac{\sqrt{3}}{6}=\boxed{(C)\frac{2\sqrt{3}}{3}}$.

 



\textbf{Solution 3}

We split rhombus $ABCD$ into two equilateral triangles, $ABD$ and $BCD$. In triangle $ABD$, the probability that a randomly selected point is closer to $B$ than either other point is $\frac{1}{3}$ *. Similarly, in triangle $BCD$, the same principle applies. Thus, the area of the region closer to $B$ than $A$, $C$, or $D$ is $\frac{1}{3} [ABD] + \frac{1}{3} [BCD]$. Since $ABD$ and $BCD$ are congruent, we have $\frac{1}{3} [ABD] + \frac{1}{3} [BCD] = \frac{2}{3} [ABD] = \frac{2}{3} \cdot \frac{s^2\sqrt{3}}{4} = \frac{2}{3} \cdot \frac{(2)^2\sqrt{3}}{4} = \boxed{\frac{2\sqrt3}{3} = C}$, and we are done.

 \*See note

 



\textbf{Solution 4}


                    \begin{figure}[H]
                        \centering
                        \includegraphics[width=0.5\textwidth]{Solution/AMC_12_12B_2011_16_Solution_Figure_2.png}
                    \end{figure}
                    Since $H$ and $E$ are halfway between $AB$ and $CB$, respectively, we know that $\overline{BH}=\overline{BE}=1$. By symmetry, $\Delta BFG$ is equilateral, so $\angle FBG=60^\circ\implies\angle EBF=\angle HBG=30^\circ$ and therefore $\Delta EBF$ and $\Delta HBG$ are 30-60-90 right triangles. Thus, $[\Delta EBF]=[\Delta BFG]=\dfrac1{2\sqrt3}$. We know that $\overline{FB}=\overline{GB}=\dfrac2{\sqrt3}$, so therefore $[\Delta BFG]=\dfrac{\sqrt3}4\left(\dfrac2{\sqrt3}\right)^2=\dfrac1{\sqrt3}$. Summing these three regions, we get $\dfrac1{2\sqrt3}+\dfrac1{2\sqrt3}+\dfrac1{\sqrt3}=\boxed{\textbf{(C)}~\dfrac{2\sqrt3}3}$. ~ Technodoggo, Asymptote diagram modified from Solution 1

 



\textbf{Solution 5}


                    \begin{figure}[H]
                        \centering
                        \includegraphics[width=0.5\textwidth]{Solution/AMC_12_12B_2011_16_Solution_Figure_3.png}
                    \end{figure}
                    To keep it simple, break rhombus $ABCD$ into two triangles, $ABD$ and $BCD$. To see the area closest to the point $B$, notice that a third of each triangle, which contains all the points nearest to $B$ in each triangle, is easily visualizable. Thus, a third of rhombus $ABCD$ must be found.

 We find the total area of rhombus $ABCD$, which we can again split into two congruent equilateral triangles with side length $2$. Using the formula of equilateral triangles and then multiplying by $\dfrac{1}{3}$:
 \[\dfrac{\sqrt{3}}{4}\cdot 2^2 \cdot 2 \cdot \dfrac{1}{3} = 2\sqrt{3} \cdot \dfrac{1}{3} = \boxed{\textbf{(C)}~\dfrac{2\sqrt{3}}{3}}\]-NSAoPS, diagram modified from Solution 1.

 \textbf{Note} There is an intuitive explanation for why the shaded region is one third of the rhombus. Consider one of the equilateral triangles; the figure is symmetrical, so one half of the shaded region should be in the triangle. Note that the perpendicular bisectors we constructed in previous solutions to deliminate $R$ must be medians, as the triangle is equilateral. Therefore, recalling that the three intersecting medians of a triangle divide it into 6 smaller triangles of equal area, we observe that the shaded region in the equilateral triangles includes precisely two of these smaller triangles. We can repeat this with the other equilateral triangle that forms the rhombus, and thus we conclude that the shaded area is one third of the rhombus's total area. 

 


	\item (\textit{AMC 12 12B 2011}) Let $f(x) = 10^{10x}, g(x) = \log_{10}\left(\frac{x}{10}\right), h_1(x) = g(f(x))$, and $h_n(x) = h_1(h_{n-1}(x))$ for integers $n \geq 2$. What is the sum of the digits of $h_{2011}(1)$?


 $\textbf{(A)}\ 16081 \qquad \textbf{(B)}\ 16089 \qquad \textbf{(C)}\ 18089 \qquad \textbf{(D)}\ 18098 \qquad \textbf{(E)}\ 18099$

 \textbf{Answer: }B



\textbf{Solution}

$g(x)=\log_{10}\left(\frac{x}{10}\right)=\log_{10}\left({x}\right) - 1$
$h_{1}(x)=g(f(x))\text{ = }g(10^{10x})=\log_{10}\left({10^{10x}}\right){ - 1 = 10x - 1}$

 Proof by induction that $h_{n}(x)\text{ = }10^n x - (1 + 10 + 10^2 + ... + 10^{n-1})$:

 For $n=1$, $h_{1}(x)=10x - 1$

 Assume $h_{n}(x)=10^n x - (1 + 10 + 10^2 + ... + 10^{n-1})$ is true for n:

 
 \begin{align*} h_{n+1}(x)&= h_{1}(h_{n}(x))\\ &=10 h_{n}(x) - 1\\ &=10 (10^n x   - (1 + 10 + 10^2 + ... + 10^{n-1})) - 1\\ &= 10^{n+1} x - (10 + 10^2 + ... + 10^{n}) - 1\\ &= 10^{n+1} x - (1 + 10 + 10^2 + ... + 10^{(n+1)-1}) \end{align*}


 Therefore, if it is true for n, then it is true for n+1; since it is also true for n = 1, it is true for all positive integers n.

 $h_{2011}(1) = 10^{2011}\times 1{ - }(1 + 10 + 10^2 + ... + 10^{2010})$, which is the 2011-digit number 8888...8889

 The sum of the digits is 8 times 2010 plus 9, or $\boxed{16089\textbf{(B)}}$




\textbf{Solution 2 (Quick, Non-Rigorous Trends)}

As before, $h_1(x)=10x-1$. Compute $h_1(x)$, $h_2(x)$, and $h_3(x)$ to yield 9, 89, and 889. Notice how this trend will repeat this trend (multiply by 10, subtract 1, repeat). As such, $h_{2011}$ is just 2010 8's followed by a nine. $2010(8)+9=\boxed{\textbf{B)}16089}$.

 


	\item (\textit{AMC 12 12B 2011}) A pyramid has a square base with side of length 1 and has lateral faces that are equilateral triangles. A cube is placed within the pyramid so that one face is on the base of the pyramid and its opposite face has all its edges on the lateral faces of the pyramid. What is the volume of this cube?


 $\textbf{(A)}\ 5\sqrt{2} - 7 \qquad \textbf{(B)}\ 7 - 4\sqrt{3} \qquad \textbf{(C)}\ \frac{2\sqrt{2}}{27} \qquad \textbf{(D)}\ \frac{\sqrt{2}}{9} \qquad \textbf{(E)}\ \frac{\sqrt{3}}{9}$

 \textbf{Answer: }A



\textbf{Solution 1}

We can use the Pythagorean Theorem to split one of the triangular faces into two 30-60-90 triangles with side lengths $\frac{1}{2}, 1$ and $\frac{\sqrt{3}}{2}$.

 Next, take a cross-section of the pyramid, forming a triangle with the top of the pyramid and the midpoints of two opposite sides of the square base.

 This triangle is isosceles with a base of 1 and two sides of length $\frac{\sqrt{3}}{2}$.

 The height of this triangle will equal the height of the pyramid. To find this height, split the triangle into two right triangles, with sides $\frac{1}{2}, \frac{\sqrt2}{2}$ and $\frac{\sqrt{3}}{2}$.

 The cube, touching all four triangular faces, will form a similar pyramid that sits on top of the cube. If the cube has side length $x$, the small pyramid has height $\frac{x\sqrt{2}}{2}$ (because the pyramids are similar).

 Thus, the height of the cube plus the height of the smaller pyramid equals the height of the larger pyramid.

 $x +\frac{x\sqrt{2}}{2} = \frac{\sqrt2}{2}$.

 $x\left(1+\frac{\sqrt{2}}{2} \right) =\frac{\sqrt{2}}{2}$
$x\left(2+\sqrt{2}\right) = \sqrt{2}$
$x = \frac{\sqrt{2}}{2+\sqrt{2}}  \cdot \frac{2-\sqrt{2}}{2-\sqrt{2}} = \frac{2\sqrt{2}-2}{4-2} = \sqrt{2}-1 =$side length of cube.

 $\left(\sqrt{2}-1\right)^3 = (\sqrt{2})^3 + 3(\sqrt{2})^2(-1) + 3(\sqrt{2})(-1)^2 + (-1)^3 = 2\sqrt{2} - 6 +3\sqrt{2} - 1 =\textbf{(A)} 5\sqrt{2} - 7$




\textbf{Solution 2}

Let the side length of the square be $s$. Let the top of the pyramid be $A$ and the square base be $BCDE$. Then, let the smaller cube meet edge $AB$ at $F$ and edge $AC$ at $G$. Let the closest vertice of the cube to $C$ on the base of the pyramid be $H$.Consider diagonal $CE$. It has length $\sqrt{2}$. Since the diagonal of the smaller cube's base is $s\sqrt{2}$, note that the distance from $C$ to $H$ is $\dfrac{\sqrt{2} - s\sqrt{2}}{2}.$ Also, note that $AG = s$ and $CG = 1-s$. We can now use the Pythagorean Theorem on triangle $CHG$, the right angle at $G$, to solve for $s$. 
 \[s^2 + (\dfrac{\sqrt{2} - s\sqrt{2}}{2})^2 = (1-s)^2\]
 \[\rightarrow s = \sqrt{2} - 1 \rightarrow \boxed{\textbf{(A)}}.\]

 \textbf{AMC 10} (For eventual redirection of one of the two pages.) This is also \#22 on the corresponding AMC 10B.

 


	\item (\textit{AMC 12 12B 2011}) A lattice point in an $xy$-coordinate system is any point $(x, y)$ where both $x$ and $y$ are integers. The graph of $y = mx + 2$ passes through no lattice point with $0 < x \leq 100$ for all $m$ such that $\frac{1}{2} < m < a$. What is the maximum possible value of $a$?


 $\textbf{(A)}\ \frac{51}{101} \qquad \textbf{(B)}\ \frac{50}{99} \qquad \textbf{(C)}\ \frac{51}{100} \qquad \textbf{(D)}\ \frac{52}{101} \qquad \textbf{(E)}\ \frac{13}{25}$

 \textbf{Answer: }B



\textbf{Solution 1}

It is very easy to see that the $+2$ in the graph does not impact whether it passes through the lattice. 

 We need to make sure that $m$ cannot be in the form of $\frac{a}{b}$ for $1\le b\le 100$. Otherwise, the graph $y= mx$ passes through the lattice point at $x = b$. We only need to worry about $\frac{a}{b}$ very close to $\frac{1}{2}$, $\frac{n+1}{2n+1}$, $\frac{n+1}{2n}$ will be the only case we need to worry about and we want the minimum of those, clearly for  $1\le b\le 100$, the smallest is $\frac{50}{99}$, so answer is $\boxed{\frac{50}{99} \textbf{(B)}}$ (In other words we are trying to find the smallest $m>\frac{1}{2}$ such that $b\le 100$.)

 



\textbf{Solution 2}

Like in the first solution, we note that the $+2$ does not affect our answer. Thus, we can ignore it and consider an equivalent problem with the line $y = mx$.

 A line with a slope of $\frac{1}{2}$ passes through (among other lattice points) the lattice point $(100,50)$. As the slope of the line increases from $\frac{1}{2}$, the first lattice point it hits is at $(99, 50)$, the slope of that line being $\frac{50}{99}$ . So the answer is $\boxed{\textbf{(B)}}$




\textbf{Solution 3}

Like in the first solution, we note that the $+2$ does not affect our answer. Thus, we can ignore it and consider an equivalent problem with the line $y = mx$.

 We will start with the first answer choice. Notice how it is impossible for a line with slope $\frac{51}{101}$ to pass through a lattice point when $0 < x \leq 100$. 

 However, if the slope is $\frac{50}{99}$, then the line will pass through the point $(99, 50)$. This sets the upper bound for $a$ as $\boxed{\textbf{(B)}}$.

 


	\item (\textit{AMC 12 12B 2011}) Triangle $ABC$ has $AB = 13, BC = 14$, and $AC = 15$. The points $D, E$, and $F$ are the midpoints of $\overline{AB}, \overline{BC}$, and $\overline{AC}$ respectively. Let $X \not= E$ be the intersection of the circumcircles of $\Delta BDE$ and $\Delta CEF$. What is $XA + XB + XC$?


 $\textbf{(A)}\ 24 \qquad \textbf{(B)}\ 14\sqrt{3} \qquad \textbf{(C)}\ \frac{195}{8} \qquad \textbf{(D)}\ \frac{129\sqrt{7}}{14} \qquad \textbf{(E)}\ \frac{69\sqrt{2}}{4}$

 \textbf{Answer: }C



\textbf{Solution 1 (Coordinates)}

Let us also consider the circumcircle of $\triangle ADF$.

 Note that if we draw the perpendicular bisector of each side, we will have the circumcenter of $\triangle ABC$ which is $P$, Also, since $m\angle ADP = m\angle AFP = 90^\circ$. $ADPF$ is cyclic, similarly, $BDPE$ and $CEPF$ are also cyclic. With this, we know that the circumcircles of $\triangle ADF$, $\triangle BDE$ and $\triangle CEF$ all intersect at $P$, so $P$ is $X$.

 The question now becomes calculating the sum of the distance from each vertex to the circumcenter. 

 We can calculate the distances with coordinate geometry. (Note that $XA = XB = XC$ because $X$ is the circumcenter.)

 Let $A = (5,12)$, $B = (0,0)$, $C = (14, 0)$, $X= (x_0, y_0)$

 Then $X$ is on the line $x = 7$ and also the line with slope $-\frac{5}{12}$ that passes through $(2.5, 6)$ (realize this is due to the fact that $XD$ is the perpendicular bisector of $AB$). 

 $y_0 = 6-\frac{45}{24} = \frac{33}{8}$

 So $X = (7, \frac{33}{8})$

 and $XA +XB+XC = 3XB = 3\sqrt{7^2 + \left(\frac{33}{8}\right)^2} = 3\times\frac{65}{8}=\boxed{\frac{195}{8}}$

 Remark: the intersection of the three circles is called a Miquel point.

 



\textbf{Solution 2 (Algebra)}

Consider an additional circumcircle on $\triangle ADF$.  After drawing the diagram, it is noticed that each triangle has side values: $7$, $\frac{15}{2}$, $\frac{13}{2}$.  Thus they are congruent, and their respective circumcircles are. 

 
Let $M$ & $N$ be $\triangle BDE$ & $\triangle CEF$'s circumcircles' respective centers. Since $\triangle BDE$ & $\triangle CEF$ are congruent, the distance $M$ & $N$ each are from $\overline{BC}$ are equal, so $\overline{MN} || \overline{BC}$. The angle between $\overline {MN}$ & $\overline{EX}$ is $90^{\circ}$, and since $\overline{MN} || \overline{BC}$, $\angle XEC$ is also $90^{\circ}$. $\triangle XEC$ is a right triangle inscribed in a circle, so $\overline{XC}$ must be the diameter of $N$. Using the same logic & reasoning, we could deduce that $XA$ & $XB$ are also circumdiameters.

 
Since the circumcircles are congruent, circumdiameters $XA$, $XB$, and $XC$ are congruent.  Therefore, the solution can be found by calculating one of these circumdiameters and multiplying it by a factor of $3$.  We can find the circumradius quite easily with the formula $\sqrt{(s)(s-a)(s-b)(s-c)} = \frac{abc}{4R}$, such that $s=\frac{a+b+c}{2}$ and $R$ is the circumradius.  Since $s = \frac{21}{2}$:

 
 \[\sqrt{(\frac{21}{2})(4)(3)(\frac{7}{2})} = \frac{\frac{15}{2}\cdot\frac{13}{2}\cdot 7}{4R}\]

 After a few algebraic manipulations:

 $\Rightarrow R=\frac{65}{16} \Rightarrow XA = XB = XC = \frac{65}{8} \Rightarrow XA + XB + XC = \boxed{\frac{195}{8}}$.

 



\textbf{Solution 3 (Dilation)}

Let $O$ be the circumcenter of $\triangle ABC,$ and $h_A$ denote the length of the altitude from $A.$ Note that a dilation centered at $B$ with ratio $\frac{1}{2}$ takes the circumcircle of $\triangle BAC$ to the circumcircle of $\triangle BDE$. It also takes the point diametrically opposite $B$ on the circumcircle of $\triangle BAC$ to $O.$ Therefore, $O$ lies on the circumcircle of $\triangle BDE.$ Similarly, it lies on the circumcircle of $\triangle CEF.$ By Pythagorean triples, $h_A=12.$ Finally, our answer is 
 \[3R=3\cdot \frac{abc}{4\{ABC\}}=3\cdot \frac{abc}{2ah_A}=3\cdot \frac{bc}{2h_A}=\boxed{\frac{195}{8}}.\]




\textbf{Solution 4 (basically Solution 1 but without coordinates)}

Since Solution 1 has already proven that the circumcenter of $\triangle ABC$ coincides with $X$, we'll go from there. Note that the radius of the circumcenter of any given triangle is $\frac{a}{2\sin{A}}$, and since $b=15$ and $\sin{B}=\frac{12}{13}$, it can be easily seen that $XA = XB = XC = \frac{65}{8}$ and therefore our answer is 
 \[3\cdot \frac{65}{8}=\boxed{\frac{195}{8}}.\]




\textbf{Solution 5}


                    \begin{figure}[H]
                        \centering
                        \includegraphics[width=0.5\textwidth]{Solution/AMC_12_12B_2011_20_Solution_Figure_1.png}
                    \end{figure}
                    

 

 Since $ED$ is a midline of $\triangle CAB,$ we have that $\triangle CED \sim \triangle CAB$ with a side length ratio of $1:2.$ 

 Consider a homothety of scale factor $2$ with on $\triangle CED$ concerning point $C$. Note that this sends $(CEDX)$ to $(ABCC')$ with $CX=XC'.$ By properties of homotheties, $C,X,$ and $C'$ are collinear. Similarly, we obtain that $BX=XB',$ with all three points collinear. Let $O$ denote the circumcenter of $\triangle ABC.$ It is well-known that $OX \perp CC'$ and analogously $OX \perp BB'.$ However, there is only one perpendicular line to $OX$ passing through $X,$, therefore, $O$ coincides with $X.$

 It follows that $AX=BX=CX=R,$ where $R$ is the circumradius of $\triangle ABC,$ and this can be computed using the formula 
 \[R=\frac{abc}{4[ABC]},\] from which we quickly obtain 
 \[R=\frac{65}{8} \implies AX+BX+CX=\boxed{\frac{195}{8}}.\]




\textbf{Solution 6 (Trigonometry)}


                    \begin{figure}[H]
                        \centering
                        \includegraphics[width=0.5\textwidth]{Solution/AMC_12_12B_2011_20_Solution_Figure_2.png}
                    \end{figure}
                    $\angle BXE = \angle BDE$, $\angle CXE = \angle CFE$, as the angles are on the same circle.

 $\triangle BDE \sim \triangle ABC$, $\triangle CFE \sim \triangle ABC$
$\angle BDE = \angle A$, $\angle CFE = \angle A$
$\angle BXE = \angle A$, $\angle CXE = \angle A$

 Therefore $\angle BXE = \angle CXE$, and $XE$ is the angle bisector of $\triangle XBC$. By the angle bisector theorem $\frac{XB}{XC} = \frac{BE}{CE} = 1$, $XB = XC$. In a similar fashion $XA = XB = XC = R$, where $R$ is the circumcircle of $\triangle ABC$.

 By the law of cosine, $\cos A = \frac{13^2 + 15^2 - 14^2}{2 \cdot 13 \cdot 15} = \frac{33}{65}$, $\sin A = \sqrt{1 - \left(\frac{33}{65}\right)^2} = \frac{56}{65}$

 By the extended law of sines, $2R = \frac{BC}{\sin A} = \frac{14}{\frac{56}{65}} = \frac{65}{4}$, $R = \frac{65}{8}$
$XA + XB + XC = 3 R = \boxed{\textbf{(C) } \frac{195}{8}}$




\textbf{Solution 7 (abwabwabwa)}

Claim, $X$ is the circumcenter of triangle $\triangle{ABC}$.

 
Proof: Note that $\triangle{BDE}$ and $\triangle{EFC}$ are congruent. Consider the centers $O_1$ and $O_2$ of $\triangle{BDE}$ and $\triangle{EFC}$, respectively. Let $B'$ be the reflection of $B$ over $O_1$, and let $C'$ be the reflection of $C$ over $O_2$. Since they form diameters, they must form right triangles $\triangle{BEB'}$ and $\triangle{CEC'}$. However, because $\triangle{BDE} \cong \triangle{EFC}$, C' and B' are the same point. Thus, one point lies on both circumcircles, so this point is $X$. But then X lies on the perpendicular bisector of $BC$, and appyling this logic to all 3 sides, $X$ must be the circumcenter.

 Memorizing that the circumradius of a $13, 14, 15$ triangle is $\frac{65}{8}$, since $XA=XB=XC=\frac{65}{8}$, $XA+XB+XC = \boxed{\textbf{(C) }\frac{195}{8}}$.

 -skibbysiggy

 


	\item (\textit{AMC 12 12B 2011}) The arithmetic mean of two distinct positive integers $x$ and $y$ is a two-digit integer. The geometric mean of $x$ and $y$ is obtained by reversing the digits of the arithmetic mean. What is $|x - y|$?


 $\textbf{(A)}\ 24 \qquad \textbf{(B)}\ 48 \qquad \textbf{(C)}\ 54 \qquad \textbf{(D)}\ 66 \qquad \textbf{(E)}\ 70$

 \textbf{Answer: }D



\textbf{Solution 1}

Answer: (D)

 $\frac{x + y}{2} = 10 a+b$ for some $1\le a\le 9$, $0\le b\le 9$.

 $\sqrt{xy} = 10 b+a$

 Squaring the first and second equations, $\frac{x^2 + 2xy + y^2}{4}=100 a^2 + 20 ab + b^2$
$xy = 100b^2 + 20ab + a^2$

 Subtracting the previous two equations, $\frac{x^2 + 2xy + y^2}{4} - xy = \frac{x^2 - 2xy + y^2}{4} = \left(\frac{x-y}{2}\right)^2 = 99 a^2 - 99 b^2 = 99(a^2 - b^2)$
$|x-y| = 2\sqrt{99(a^2 - b^2)}=6\sqrt{11(a^2 - b^2)}$

 Note that for x-y to be an integer, $(a^2 - b^2)$ has to be $11n$ for some perfect square $n$. Since $a$ is at most $9$, $n = 1$ or $4$

 If $n = 1$, $|x-y| = 66$, if $n = 4$, $|x-y| = 132$. In AMC, we are done. Otherwise, we need to show that $n=4$, or $a^2 -b^2 = 44$ is impossible.

 $(a-b)(a+b) = 44$ -> $a-b = 1$, or $2$ or $4$ and $a+b = 44$, $22$, $11$ respectively. And since $a+b \le 18$, $a+b = 11$, $a-b = 4$, but there is no integer solution for $a$, $b$.

 \textbf{Short Cut} We can arrive at $|x-y| = 6\sqrt{11(a^2 - b^2)}$ using the method above. Because we know that $|x-y|$ is an integer, it must be a multiple of 6 and 11. Hence the answer is $66.$

 In addition:Note that $11n$ with $n = 1$ may be obtained with $a = 6$ and $b = 5$ as $a^2 - b^2 = 36 - 25 = 11$.

 \textbf{Sidenote} It is easy to see that $(a,b)=(6,5)$ is the only solution. This yields $(x,y)=(98,32)$. Their arithmetic mean is $65$ and their geometric mean is $56$.

 



\textbf{Solution 2}

Let $(x+y)/2 = 10a + b$ and $\sqrt{xy} = 10b + a$. By AM-GM we know that $a \ge b$. Squaring and multiplying by 4 on the first equation we get $x^2 + y^2 + 2xy = 400a^2 + 4b^2 + 80ab$. Squaring and multiplying the second equation by 4 we get $4xy = 400b^2 + 4a^2 + 80ab$. Subtracting we get $(x-y)^2 = 396(a^2 - b^2)$. Note that $396 = 2^2 \cdot 3^2 \cdot 11$. So to make it a perfect square $a^2 - b^2 = 11$. From \href{https://artofproblemsolving.com/wiki/index.php?title=Difference\_of\_squares}{difference of squares}, we see that $a = 6$ and $b = 5$. So the answer is $3 \cdot 2 \cdot 11 = 66$. ~coolmath\_2018

 


 



\textbf{Solution 3 (whee!)}

Let the 2 digit number be $10a+b$. We know $x+y=20a+2b$ and $\sqrt{xy} = 10b+a$.

 
Thus. $x^2+2xy+y^2 = 400a^2+80ab+4b^2$ and $xy = 100b^2+20ab+a^2$. Notice that $|x-y| = \sqrt{(x-y)^2}$. We know that $(x-y)^2 = (x+y)^2-4xy = 396a^2-396b^2$. Then, $|x-y| = 6\sqrt{11}\cdot \sqrt{b^2-a^2}$. Clearly, $|x-y|$ must be a multiple of 11, so the only possible answer is $\boxed{\textbf{(D) }66}$.

 -skibbysiggy

 


	\item (\textit{AMC 12 12B 2011}) Let $T_1$ be a triangle with side lengths $2011, 2012$, and $2013$. For $n \geq 1$, if $T_n = \Delta ABC$ and $D, E$, and $F$ are the points of tangency of the incircle of $\Delta ABC$ to the sides $AB, BC$, and $AC$, respectively, then $T_{n+1}$ is a triangle with side lengths $AD, BE$, and $CF$, if it exists. What is the perimeter of the last triangle in the sequence $\left(T_n\right)$?


 $\textbf{(A)}\ \frac{1509}{8} \qquad \textbf{(B)}\  \frac{1509}{32} \qquad \textbf{(C)}\  \frac{1509}{64} \qquad \textbf{(D)}\  \frac{1509}{128} \qquad \textbf{(E)}\  \frac{1509}{256}$

 \textbf{Answer: }D



\textbf{Solution}

Answer: (D)

 Let $AB = c$, $BC = a$, and $AC = b$

 Then $AD = AF$, $BE = BD$ and $CF = CE$

 Then $a = BE + CF$, $b = AD + CF$, $c = AD + BE$

 Hence:

 \begin{verbatim}
$AD = AF = \frac{b + c - a}{2}$$BE = BD = \frac{a + c - b}{2}$$CF = CE = \frac{a + b - c}{2}$
\end{verbatim}

Note that $a + 1 = b$ and $a - 1 = c$ for $n = 1$, I claim that it is true for all $n$, assume for induction that it is true for some $n$, then

 \begin{verbatim}
$AD = AF = \frac{a}{2}$$BE = BD = \frac{a - 2}{2} = AD - 1$$CF = CE = \frac{a + 2}{2} = AD + 1$
\end{verbatim}

Furthermore, the average for the sides is decreased by a factor of 2 each time.

 So $T_n$ is a triangle with side length $\frac{2012}{2^{n- 1}} - 1$, $\frac{2012}{2^{n-1}}$, $\frac{2012}{2^{n-1}} + 1$

 and the perimeter of such $T_n$ is $\frac{(3)(2012)}{2^{n-1}}$

 
Now we need to find when $T_n$ fails the triangle inequality. So we need to find the last $n$ such that $\frac{2012}{2^{n-1}} > 2$
\begin{verbatim}
$\frac{2012}{2^{n-1}} > 2$$2012 > 2^n$$n \le 10$
\end{verbatim}

For $n = 10$, perimeter is $\frac{(3)(2012)}{2^{9}} = \frac{1509}{2^7} = \frac{1509}{128}$




\textbf{Solution 2}

This problem is a PoP(Power of Point) Bash exercise. What I did to solve this problem was look at recursive perimeters. The first ever perimeter is $P_1 = 2011 + 2012 + 2013$. Next, by PoP, inscribing the circle gives us three new lengths, namely $AD, BE, CF$. Denote $AD$ as $x_1$ and homogeneously the others. Then, $x_1 + x_2 = 2011, x_2 + x_3 = 2012$, and $x_1 + x_3 = 2013$. If we add all these equations up and divide by $2$, we get $x_1 + x_2 + x_3 = \frac{2011 + 2012 + 2013}{2}$. Writing this with $P_1 = 2011 + 2012 + 2013$, we get that our new perimeter, $P_2$, is indeed equal to $\frac{P_1}{2}$. Similarly, by the same concept, we get that $P_3 = \frac{P_1}{4}$ and the pattern keeps going. In general, I found that for each new perimeter $P_n$, $P_n = \frac{P_1}{2^{n-1}}$. Now, substituting in the numerical value of $P_1$, we get that $P_n = \frac{2011 + 2012 + 2013}{2^{n-1}}$. If you keep dividing the numerator and denominator by 2, I got that: $P_n = \frac{1509}{2^{n-3}}$. This representation of a new perimeter in terms of $n$ looks very similar to the option choices, so we're on the right path. Now, all we need to do is find out the last value of $n$ when the sequence stops working. By the Triangular Inequality, we may be able to finish this off. Now, I won't go into depths here, but listing out terms and using the Triangular Inequality gives us:First Sequence:$x_1 < 2012$
$x_2 < 2013$
$x_3 < 2011$.The ordered triple for this set of $(x_1, x_2, x_3)$ is $(1005, 1006, 1007)$ if you solve the PoP system of equations we got earlier. If you keep listing out terms, we can come up with the general form, and that is:$x_1 < \frac{503}{2^{n-4}}$
$x_2 < \frac{503}{2^{n-4}} + 1$
$x_3 < \frac{503}{2^{n-4}} - 1$.The ordered triple for this set of $(x_1, x_2, x_3)$ is $(\frac{503}{2^{n-4}}, \frac{503}{2^{n-4}} - 1, \frac{503}{2^{n-4}} + 1)$. Now, if we plug in these values of $x_1, x_2$, and $x_3$ into the inequalities, we see that the first two are always satisfied, but the last one is only satisfied when:$2^{n-2} < 503$(You will get this if you simplify the last inequality). The last $n$ for this to be satisfied is when $n = 10$. If we go up to our general representation of $P_n$, we see that $P_n = \frac{1509}{2^{n-3}}$. Plugging in $10$ because this is the last $n$(and also the last triangle), we our final answer of $\frac{1509}{2^{7}} = \boxed{\frac{1509}{128}}$ or $\boxed{D}$.

 


	\item (\textit{AMC 12 12B 2011}) A bug travels in the coordinate plane, moving only along the lines that are parallel to the $x$-axis or $y$-axis. Let $A = (-3, 2)$ and $B = (3, -2)$. Consider all possible paths of the bug from $A$ to $B$ of length at most $20$. How many points with integer coordinates lie on at least one of these paths?


 $\textbf{(A)}\ 161 \qquad \textbf{(B)}\ 185 \qquad \textbf{(C)}\  195 \qquad \textbf{(D)}\  227 \qquad \textbf{(E)}\  255$

 \textbf{Answer: }C



\textbf{Solution}

We declare a point $(x, y)$ to make up for the extra steps that the bug has to move. If the point $(x, y)$ satisfies the property that $|x - 3| + |y + 2| + |x + 3| + |y - 2| \le 20$, then it is in the desirable range because $|x - 3| + |y + 2|$ is the length of the shortest path from $(x,y)$ to $(3, -2)$ and $|x + 3| + |y - 2|$ is the length of the shortest path from $(x,y)$ to $(-3, 2)$.

 
If $-3\le x \le 3$, then $-7\le y \le 7$ satisfy the property. there are $15 \times 7 = 105$ lattice points here.

 else let $3< x \le 8$ (and for $-8 \le x < -3$ because it is symmetrical) We set 8 as the upper bound for x because the shortest distance from $(-3, 2)$ to $(x, y)$ added to the shortest distance from $(x, y)$ to $(3, -2)$ is $|x - 3| + |y + 2| + |x + 3| + |y - 2|$. Since the minimum value for the difference between the y-coordinates is at $y = 0$, we get $2x + 4 = 20$ or $-2x + 4 = 20$. Thus, the upper and lower bounds for $x$ are $8$ and $-8$, respectively.

 Now we test each value for x satisfying $3< x \le 8$ and double the result because of symmetry.

 For $x = 4$, the possibles values of y are such that $|2y| \le 12$ for a total of $13$ lattice points,

 for $x = 5$, the possibles values of y are such that $|2y| \le 10$ for a total of $11$ lattice points,

 for $x = 6$, the possibles values of y are such that $|2y| \le 8$ for a total of $9$ lattice points,

 for $x = 7$, the possibles values of y are such that $|2y| \le 6$ for a total of $7$ lattice points,

 for $x = 8$, the possibles values of y are such that $|2y| \le 4$ for a total of $5$ lattice points,


 Hence, there are a total of $105 + 2 ( 13 + 11 + 9 + 7 + 5) = \boxed{(C) 195}$ lattice points.

 One may also obtain the result by using Pick's Theorem(how?).

 $i = a - \frac{b}{2} - 1$ (Suggestion)

 



\textbf{Solution 2}


                    \begin{figure}[H]
                        \centering
                        \includegraphics[width=0.5\textwidth]{Solution/AMC_12_12B_2011_23_Solution_Figure_1.png}
                    \end{figure}
                    (Diagram for Solution 2 by Technodoggo also don't ask about the patterns, I just figured out how to use the patterns module)

 Notice that the bug is basically moving from A to B (length 10) but optionally going on a detour anywhere along its route. 

 Specifically, the detour would be of total length 10, 5 to some point and then going back by retracing its path (also length 5). Just to simplify things we can observe that the coordinates don't matter here and all we need to remember is that A and B are the diagonally opposite vertices of a 4 by 6 rectangle (5 by 7 lattice points). The bug can start the "detour" from any point on or inside the rectangle. Notice that the bug can go 5 steps in the y direction, 5 steps in the x direction, or anything in between, so the points covered by possible detours from any point would look like a rhombus or square rotated 45 degrees (with centre at a point on or inside the rectangle). Drawing this out we would get an octagon. 

 Finding the final answer is then easy, for this solution I will slice the octagon into 4 rectangles (2 of which are squares) and 4 isosceles triangles. There are $(4+1)\cdot (6+1) = 35$ points on or inside the original triangle, $5\cdot 7$ points covered by the rectangles above and below the original one and $5\cdot 5$ points for the squares to the right and left of the original triangle. Lastly each of the four isosceles triangles cover $4+3+2+1 = 10$ points. (Notice that although the length of the detour is 5, the points on the edge of the triangles were already counted). 

 Adding these up, we get $3\cdot 35 + 2\cdot 25 + 4\cdot 10 = 195 => \boxed{(C)}$

 Another way to get to the octagon would be to note that aside from the 4x6/5x7 box of the bare absolutely necessary route (the checkered region), we can go on a 5-step detour anywhere (as above); we need only examine how far we can get through the edges (the border), because that takes us the farthest. Starting from any point on the border, the farthest that it can go is 5 units in some cardinal direction, which gives us the tiled region in the diagram above. Finally, from our corners, we can also go 4-1, 3-2, 2,3, or 1-4, which gives us the rest of the region enclosed in blue. This gives us our octagon. 

 tl;dr the interior of the octagon in blue is a result of taking all the points that can be reached from the boundary of the box as described in no more than 5 moves. We can then proceed to find the area in any simple way. (Extra way & diagram by Technodoggo)

 


 



\textbf{Solution 3 (think about it a lil bit)}

We start with reaching $B$ from $A$ first. It only takes $10$ moves, a sequence of moves down and towards the right. This is a $5$ by $7$ grid of squares. For the remaining $10$ moves, we can move $5$ moves out, and then spend the other $5$ moves "undoing" these $5$ moves. For example, we can move $4$ left and $1$ up at the start, and of the remaining $15$ moves, there must be $4$ "downs" and $6$ "rights", along with $5$ moves that undo the original $5$ at the start.

 
Drawing it on a grid, we will have one center rectangle, 4 rectangles attached to it (1 on each side), and 1 triangle in each corner. The total number of points is $35+35+35+25+25+10+10+10+10 = \boxed{\textbf{(C) }195}$.

 -skibbysiggy

 


	\item (\textit{AMC 12 12B 2011}) Let $P(z) = z^8 + \left(4\sqrt{3} + 6\right)z^4 - \left(4\sqrt{3} + 7\right)$. What is the minimum perimeter among all the $8$-sided polygons in the complex plane whose vertices are precisely the zeros of $P(z)$?


 $\textbf{(A)}\ 4\sqrt{3} + 4 \qquad \textbf{(B)}\ 8\sqrt{2} \qquad \textbf{(C)}\  3\sqrt{2} + 3\sqrt{6} \qquad \textbf{(D)}\  4\sqrt{2} + 4\sqrt{3} \qquad \textbf{(E)}\  4\sqrt{3} + 6$

 \textbf{Answer: }B



\textbf{Solution}

Answer: (B)

 First of all, we need to find all $z$  such that $P(z) = 0$
$P(z) = \left(z^4 - 1\right)\left(z^4 + \left(4\sqrt{3} + 7\right)\right)$

 So $z^4 = 1 \implies z =  e^{i\frac{n\pi}{2}}$

 or $z^4 = - \left(4\sqrt{3} + 7\right)$
$z^2 = \pm i \sqrt{4\sqrt{3} + 7} = e^{i\frac{(2n+1)\pi}{2}} \left(\sqrt{3} + 2\right)$
$z = e^{i\frac{(2n+1)\pi}{4}} \sqrt{\sqrt{3} + 2} = e^{i\frac{(2n+1)\pi}{4}} \left(\frac{\sqrt{3}}{\sqrt{2}} + \frac{1}{\sqrt{2}}\right)$

 


 Now we have a solution at $\frac{n\pi}{4}$ if we look at them in polar coordinate, further more, the 8-gon is symmetric (it is an equilateral octagon) . So we only need to find the side length of one and multiply by $8$.

 So answer $= 8 \times$ distance from $1$ to $\left(\frac{\sqrt{3}}{2} + \frac{1}{2}\right)\left(1 + i\right)$

 Side length $= \sqrt{\left(\frac{\sqrt{3}}{2} - \frac{1}{2}\right)^2 + \left(\frac{\sqrt{3}}{2} + \frac{1}{2}\right)^2} = \sqrt{2\left(\frac{3}{4} + \frac{1}{4}\right)} = \sqrt{2}$

 Hence, answer is $8\sqrt{2}$.

 



\textbf{Solution 2}

Use the law of cosines.We make $a$ the distance.Now, since the angle does not change the distance from the origin, we can just use the distance. $a^2 = (\frac{\sqrt{3}}{\sqrt{2}} + \frac{1}{\sqrt{2}})^2 + 1^2 -2 \times \Big( \frac{\sqrt{3}}{\sqrt{2}} + \frac{1}{\sqrt{2}} \Big)\times 1 \times \cos \frac{\pi}{4}$, which simplifies to $a^2= 2 + \sqrt3 +1 - 1 - \sqrt3$, or $a^2=2$, or $a=\sqrt2$.Multiply the answer by 8 to get $\boxed{ (B) 8\sqrt2}$



	\item (\textit{AMC 12 12B 2011}) For every $m$ and $k$ integers with $k$ odd, denote by $\left[\frac{m}{k}\right]$ the integer closest to $\frac{m}{k}$. For every odd integer $k$, let $P(k)$ be the probability that


 
 \[\left[\frac{n}{k}\right] + \left[\frac{100 - n}{k}\right] = \left[\frac{100}{k}\right]\]


 for an integer $n$ randomly chosen from the interval $1 \leq n \leq 99!$. What is the minimum possible value of $P(k)$ over the odd integers $k$ in the interval $1 \leq k \leq 99$?


 $\textbf{(A)}\ \frac{1}{2} \qquad \textbf{(B)}\ \frac{50}{99} \qquad \textbf{(C)}\ \frac{44}{87} \qquad \textbf{(D)}\  \frac{34}{67} \qquad \textbf{(E)}\  \frac{7}{13}$

 \textbf{Answer: }D



\textbf{Solution 1}

Answer: $(D)  \frac{34}{67}$

 
First of all, you have to realize that 

 if $\left[\frac{n}{k}\right] + \left[\frac{100 - n}{k}\right] = \left[\frac{100}{k}\right]$

 then $\left[\frac{n - k}{k}\right] + \left[\frac{100 - (n - k)}{k}\right] = \left[\frac{100}{k}\right]$

 So, we can consider what happen in $1\le n \le k$ and it will repeat. Also since range of $n$ is $1$ to $99!$, it is always a multiple of $k$. So we can just consider $P(k)$ for $1\le n \le k$.

 
Let $\text{fpart}(x)$ be the fractional part function

 This is an AMC exam, so use the given choices wisely. With the given choices, and the previous explanation, we only need to consider $k = 99$, $87$, $67$, $13$. $1\le n \le k$

 
For $k > \frac{200}{3}$, $\left[\frac{100}{k}\right] = 1$. 3 of the $k$ that we should consider land in here.

 For $n < \frac{k}{2}$, $\left[\frac{n}{k}\right] = 0$, then we need $\left[\frac{100 - n}{k}\right] = 1$

 else for $\frac{k}{2}< n < k$, $\left[\frac{n}{k}\right] = 1$, then we need $\left[\frac{100 - n}{k}\right] = 0$


 For $n < \frac{k}{2}$, $\left[\frac{100 - n}{k}\right] = \left[\frac{100}{k} - \frac{n}{k}\right]= 1$

 So, for the condition to be true, $100 - n > \frac{k}{2}$ .  ( $k > \frac{200}{3}$, no worry for the rounding to be $> 1$)

 $100 > k > \frac{k}{2} + n$, so this is always true.

 For $\frac{k}{2}< n < k$, $\left[\frac{100 - n}{k}\right] = 0$, so we want $100 - n < \frac{k}{2}$, or $100 < \frac{k}{2} + n$
$100 <\frac{k}{2} + n <  \frac{3k}{2}$

 For k = 67, $67 > n > 100 - \frac{67}{2} = 66.5$

 For k = 69, $69 > n > 100 - \frac{69}{2} = 67.5$

 etc.

 
We can clearly see that for this case, $k = 67$ has the minimum $P(k)$, which is $\frac{34}{67}$. Also, $\frac{7}{13} > \frac{34}{67}$ .

 So for AMC purpose, answer is $\boxed{\textbf{(D) }\frac{34}{67}}$.

 \textbf{Proof:} Notice that for these integers $99,87,67$:

 
$0\rightarrow 49,50,51\rightarrow 98$
$100\rightarrow 51,50,49\rightarrow 2$
$P=\frac{98}{99}$
$0\rightarrow 43,44\rightarrow 56,57\rightarrow 86$
$87\rightarrow 57,56\rightarrow 44,43\rightarrow 14$
$P=\frac{74}{87}$
$0\rightarrow 33,34\rightarrow 66$
$100\rightarrow 67,66\rightarrow 34$
$P=\frac{34}{67}$

 
That the probability is $\frac{2k-100}{k}=2-\frac{100}{k}$. Even for $k=13$, $P(13)=\frac{9}{13}=\frac{100}{13}-7$. And $P(11)=\frac{10}{11}=10-\frac{100}{11}$.

 Perhaps the probability for a given $k$ is $\left\lceil{\frac{100}{k}}\right\rceil-\frac{100}{k}$ if $\left[\frac{100}{k}\right]=\left\lfloor{\frac{100}{k}}\right\rfloor$ and $\frac{100}{k}-\left\lfloor{\frac{100}{k}}\right\rfloor$ if $\left[\frac{100}{k}\right]=\left\lceil{\frac{100}{k}}\right\rceil$.

 So $P>\frac{1}{2}$ and $P_\text{min}=\frac{k_\text{min}+1}{2k_\text{min}}=\frac{101}{201}$. Because $201=3\cdot 67\mid 99!$  !

 


 


 Now, let's say we are not given any answer, we need to consider $k < \frac{200}{3}$.

 I claim that $P(k) \ge \frac{1}{2} + \frac{1}{2k}$

 
If $\left[\frac{100}{k}\right]$ got round down, then $1 \le n \le \frac{k}{2}$ all satisfy the condition along with $n = k$

 because if $\text{fpart}\left(\frac{100}{k}\right) < \frac{1}{2}$ and $\text{fpart} \left(\frac{n}{k}\right) < \frac{1}{2}$, so must $\text{fpart} \left(\frac{100 - n}{k}\right) < \frac{1}{2}$

 and for $n = k$, it is the same as $n = 0$.

 , which makes 

 $P(k) \ge \frac{1}{2} + \frac{1}{2k}$.

 
If $\left[\frac{100}{k}\right]$ got round up, then $\frac{k}{2} \le n \le k$ all satisfy the condition along with $n = 1$

 
because if $\text{fpart}\left(\frac{100}{k}\right) > \frac{1}{2}$ and $\text{fpart} \left(\frac{n}{k}\right) > \frac{1}{2}$

 Case 1)$\text{fpart} \left(\frac{100 - n}{k}\right) < \frac{1}{2}$ 

 -> $\text{fpart}\left(\frac{100}{k}\right) = \text{fpart} \left(\frac{n}{k}\right) +\text{fpart} \left(\frac{100 - n}{k}\right)$ 

 Case 2)

 $\text{fpart} \left(\frac{100 - n}{k}\right) > \frac{1}{2}$ 

 -> $\text{fpart}\left(\frac{100}{k}\right) + 1 = \text{fpart} \left(\frac{n}{k}\right) +\text{fpart} \left(\frac{100 - n}{k}\right)$ 

 
and for $n = 1$, since $k$ is odd, $\left[\frac{99}{k}\right] \neq \left[\frac{100}{k}\right]$ 

 -> $99.5 = k (p + .5)$  -> $199 = k (2p + 1)$, and $199$ is prime so $k = 1$ or $k =199$, which is not in this set

 , which makes 

 $P(k) \ge \frac{1}{2} + \frac{1}{2k}$.

 
Now the only case without rounding, $k = 1$. It must be true.

 



\textbf{Solution 2}

It suffices to consider $0\le n \le k-1.$ Now for each of these $n,$ let $f(n)=\left[\frac{n}{k}\right], g(n)=\left[\frac{100-n}{k}\right]-\left[\frac{100}{k}\right].$ If we let $k=67,$ then the following graphs result for $f$ and $g.$
$f:$ 
                    \begin{figure}[H]
                        \centering
                        \includegraphics[width=0.5\textwidth]{Solution/AMC_12_12B_2011_25_Solution_Figure_1.png}
                    \end{figure}
                    
$g:$
                    \begin{figure}[H]
                        \centering
                        \includegraphics[width=0.5\textwidth]{Solution/AMC_12_12B_2011_25_Solution_Figure_2.png}
                    \end{figure}
                    

 Our probability is the number of $0\le i \le k-1$ such that $f(i)+g(i)=0$ over $k.$ Of course, this always holds for $i=0.$ If we let $k$ vary, then the graph of $f$ is always very similar to what it looks like above (groups of $\frac{k+1}{2},\frac{k-1}{2}$ dots). However, the graph of $g$ can vary greatly. In the above diagram, $g(i)=0$ for all $i,$ while it is possible for $g(i)=-1$ for all $i\neq 0.$ In order to minimize the number of $i$ which satisfy $f(i)+g(i)=0,$ we either want $g(i)=0$ for $0\le i<k,$ or $g(i)=-1$ for $1\le i<k.$ This way, we see that at least half of the numbers from $1$ to $k-1$ satisfy the given equation. So, our desired probability is at least $\frac{k+1}{2k}.$ As shown by the diagram above, the probability is $\frac{34}{67}$ for $k=67.$ Clearly no better solutions can exist when $k<67.$ On the other hand, for $k>67$ $87$ and $99$ do not yield better probabilities. Therefore, our answer is $\boxed{\frac{34}{67}}.$



\end{enumerate}